\documentclass[12pt,a4paper]{article}
\usepackage[utf8]{inputenc}
\usepackage[margin=1in]{geometry}
\usepackage{graphicx}
\usepackage{hyperref}
\usepackage{listings}
\usepackage{xcolor}
\usepackage{tcolorbox}
\usepackage{enumitem}
\usepackage{fancyhdr}
\usepackage{titlesec}

% Colors
\definecolor{codegreen}{rgb}{0,0.6,0}
\definecolor{codegray}{rgb}{0.5,0.5,0.5}
\definecolor{codepurple}{rgb}{0.58,0,0.82}
\definecolor{backcolour}{rgb}{0.95,0.95,0.92}
\definecolor{cmdcolor}{rgb}{0.2,0.2,0.2}

% Code listing style
\lstdefinestyle{mystyle}{
    backgroundcolor=\color{backcolour},   
    commentstyle=\color{codegreen},
    keywordstyle=\color{codepurple},
    numberstyle=\tiny\color{codegray},
    stringstyle=\color{codegreen},
    basicstyle=\ttfamily\small,
    breakatwhitespace=false,         
    breaklines=true,                 
    captionpos=b,                    
    keepspaces=true,                 
    showspaces=false,                
    showstringspaces=false,
    showtabs=false,                  
    tabsize=2,
    frame=single
}
\lstset{style=mystyle}

% Header/Footer
\pagestyle{fancy}
\fancyhf{}
\rhead{SWE645 - Assignment 2}
\lhead{Divya Soni}
\rfoot{Page \thepage}

% Title formatting
\titleformat{\section}{\Large\bfseries\color{blue!70!black}}{\thesection}{1em}{}
\titleformat{\subsection}{\large\bfseries\color{blue!50!black}}{\thesubsection}{1em}{}

\begin{document}

% Title Page
\begin{titlepage}
    \centering
    \vspace*{2cm}
    {\Huge\bfseries SWE645 - Assignment 2\\[0.5cm]}
    {\Large Docker, Kubernetes \& CI/CD Pipeline\\[1cm]}
    {\large Step-by-Step Installation and Setup Guide\\[2cm]}
    
    \vfill
    
    {\Large\bfseries Divya Soni}\\[0.5cm]
    {\large SWE645}\\[0.5cm]
    {\large George Mason University}\\[1cm]
    {\large \today}
    
    \vfill
\end{titlepage}

\tableofcontents
\newpage

% ============================================
\section{Project Overview}
% ============================================

This document provides step-by-step instructions for completing SWE645 Assignment 2, which involves:

\begin{itemize}
    \item Containerizing the HW1 web application using Docker
    \item Deploying to Kubernetes using Rancher on AWS
    \item Setting up a CI/CD pipeline using Jenkins
\end{itemize}

\subsection{Architecture Diagram}

\begin{center}
\begin{tcolorbox}[width=0.9\textwidth, colback=gray!10, colframe=gray!50]
\centering
\textbf{GitHub Repository} $\rightarrow$ \textbf{Jenkins (CI/CD)} $\rightarrow$ \textbf{Docker Hub} $\rightarrow$ \textbf{Kubernetes (Rancher/AWS)}
\end{tcolorbox}
\end{center}

\subsection{Prerequisites}

\begin{itemize}
    \item AWS Account (AWS Academy)
    \item Docker Hub account (\url{https://hub.docker.com})
    \item GitHub account (\url{https://github.com})
    \item Docker Desktop installed
    \item Git installed
\end{itemize}

\newpage

% ============================================
\section{Part 1: Docker Setup}
% ============================================

\subsection{Step 1.1: Install Docker Desktop}

\begin{enumerate}
    \item Go to: \url{https://www.docker.com/products/docker-desktop}
    \item Download \textbf{Docker Desktop for Windows}
    \item Run the installer
    \item \textbf{Restart your computer}
    \item Open Docker Desktop and wait for it to start
\end{enumerate}

\textbf{Verify installation} - Open PowerShell:
\begin{lstlisting}[language=bash]
docker --version
\end{lstlisting}

Expected output: \texttt{Docker version 24.x.x}

\subsection{Step 1.2: Create Docker Hub Account}

\begin{enumerate}
    \item Go to: \url{https://hub.docker.com}
    \item Click \textbf{Sign Up}
    \item Create account and remember your username
    \item Verify your email
\end{enumerate}

\subsection{Step 1.3: Build Docker Image}

Open \textbf{PowerShell} and navigate to your project folder:

\begin{lstlisting}[language=bash]
cd C:\Users\lenovo\Desktop\ass1-645
\end{lstlisting}

Build the Docker image (replace \texttt{YOUR\_DOCKERHUB\_USERNAME}):

\begin{lstlisting}[language=bash]
docker build -t YOUR_DOCKERHUB_USERNAME/studentsurvey645:latest .
\end{lstlisting}

\subsection{Step 1.4: Test Docker Image Locally}

\begin{lstlisting}[language=bash]
docker run -d -p 8080:80 YOUR_DOCKERHUB_USERNAME/studentsurvey645:latest
\end{lstlisting}

Open browser and go to: \textbf{http://localhost:8080}

You should see your website running!

\subsection{Step 1.5: Push Image to Docker Hub}

Login to Docker Hub:
\begin{lstlisting}[language=bash]
docker login
\end{lstlisting}

Enter your Docker Hub username and password when prompted.

Push the image:
\begin{lstlisting}[language=bash]
docker push YOUR_DOCKERHUB_USERNAME/studentsurvey645:latest
\end{lstlisting}

\textbf{Verify:} Go to \url{https://hub.docker.com} and check your repository exists.

\newpage

% ============================================
\section{Part 2: Create GitHub Repository}
% ============================================

\subsection{Step 2.1: Create Repository on GitHub}

\begin{enumerate}
    \item Go to: \url{https://github.com}
    \item Click \textbf{+} $\rightarrow$ \textbf{New repository}
    \item Repository name: \texttt{studentsurvey645}
    \item Make it \textbf{Public}
    \item Click \textbf{Create repository}
\end{enumerate}

\subsection{Step 2.2: Push Code to GitHub}

Open \textbf{PowerShell}:

\begin{lstlisting}[language=bash]
cd C:\Users\lenovo\Desktop\ass1-645

git init

git add .

git commit -m "Initial commit - SWE645 HW2"

git branch -M main

git remote add origin https://github.com/YOUR_GITHUB_USERNAME/studentsurvey645.git

git push -u origin main
\end{lstlisting}

Replace \texttt{YOUR\_GITHUB\_USERNAME} with your actual GitHub username.

\newpage

% ============================================
\section{Part 3: Set Up Rancher on AWS}
% ============================================

\subsection{Step 3.1: Launch EC2 Instance for Rancher}

\begin{enumerate}
    \item Go to \textbf{AWS Academy} $\rightarrow$ \textbf{AWS Console} $\rightarrow$ \textbf{EC2}
    \item Click \textbf{Launch Instance}
    \item Configure the instance:
\end{enumerate}

\begin{center}
\begin{tabular}{|l|l|}
\hline
\textbf{Setting} & \textbf{Value} \\
\hline
Name & Rancher-Server \\
AMI & Ubuntu Server 22.04 LTS \\
Instance Type & t2.medium (NOT t2.micro) \\
Key Pair & Create new: rancher-key \\
\hline
\end{tabular}
\end{center}

\textbf{Security Group Rules:}

\begin{center}
\begin{tabular}{|l|l|l|}
\hline
\textbf{Type} & \textbf{Port} & \textbf{Source} \\
\hline
SSH & 22 & My IP \\
HTTP & 80 & Anywhere (0.0.0.0/0) \\
HTTPS & 443 & Anywhere (0.0.0.0/0) \\
Custom TCP & 6443 & Anywhere (0.0.0.0/0) \\
\hline
\end{tabular}
\end{center}

\begin{enumerate}[resume]
    \item Click \textbf{Launch Instance}
    \item Wait for \textbf{Running} status
    \item Copy the \textbf{Public IP Address}
\end{enumerate}

\subsection{Step 3.2: Connect to Rancher EC2}

Open \textbf{PowerShell}:

\begin{lstlisting}[language=bash]
cd C:\Users\lenovo\Downloads

ssh -i "rancher-key.pem" ubuntu@YOUR-RANCHER-PUBLIC-IP
\end{lstlisting}

Type \texttt{yes} when asked about fingerprint.

\subsection{Step 3.3: Install Docker on Rancher EC2}

Run these commands one by one:

\begin{lstlisting}[language=bash]
sudo apt-get update

sudo apt install -y docker.io

sudo systemctl start docker

sudo systemctl enable docker

docker --version
\end{lstlisting}

\subsection{Step 3.4: Install Rancher}

\begin{lstlisting}[language=bash]
sudo docker run --privileged -d --restart=unless-stopped \
    -p 80:80 -p 443:443 rancher/rancher
\end{lstlisting}

\textbf{Wait 2-3 minutes}, then open browser:

\begin{center}
\texttt{https://YOUR-RANCHER-PUBLIC-IP}
\end{center}

Accept the security warning (click Advanced $\rightarrow$ Proceed).

\subsection{Step 3.5: Get Rancher Bootstrap Password}

In SSH terminal:

\begin{lstlisting}[language=bash]
sudo docker logs $(sudo docker ps -q) 2>&1 | grep "Bootstrap Password"
\end{lstlisting}

Copy the password and paste it in the browser. Set a new admin password.

\subsection{Step 3.6: Create AWS IAM User for Rancher}

\begin{enumerate}
    \item Go to \textbf{AWS Console} $\rightarrow$ \textbf{IAM}
    \item Click \textbf{Users} $\rightarrow$ \textbf{Add user}
    \item User name: \texttt{rancher-user}
    \item Check \textbf{Programmatic access}
    \item Click \textbf{Next: Permissions}
    \item Click \textbf{Attach existing policies directly}
    \item Check \textbf{AdministratorAccess}
    \item Click through to \textbf{Create user}
    \item \textbf{SAVE the Access Key ID and Secret Access Key!}
\end{enumerate}

\subsection{Step 3.7: Create Kubernetes Cluster in Rancher}

\begin{enumerate}
    \item In Rancher UI, click \textbf{Create} (or \textbf{Cluster Management} $\rightarrow$ \textbf{Create})
    \item Select \textbf{Amazon EC2}
    \item Enter cluster details:
    \begin{itemize}
        \item Cluster Name: \texttt{studentsurvey-cluster}
        \item Access Key: (from Step 3.6)
        \item Secret Key: (from Step 3.6)
        \item Region: \texttt{us-east-1}
    \end{itemize}
    \item Click \textbf{Add Node Template}:
    \begin{itemize}
        \item Instance Type: \texttt{t2.medium}
        \item Give template a name
        \item Click \textbf{Create}
    \end{itemize}
    \item Configure Node Pools:
    \begin{itemize}
        \item etcd: 1 node
        \item Control Plane: 1 node
        \item Worker: 2 nodes
    \end{itemize}
    \item Click \textbf{Create}
\end{enumerate}

\textbf{Wait 10-15 minutes} for the cluster to be ready!

\subsection{Step 3.8: Get Kubeconfig File}

\begin{enumerate}
    \item In Rancher, click on your cluster
    \item Click \textbf{Kubeconfig File} (top right button)
    \item Copy all content
    \item On your local machine, create folder and file:
\end{enumerate}

\begin{lstlisting}[language=bash]
# Create .kube folder in your home directory
mkdir C:\Users\lenovo\.kube

# Create config file and paste the kubeconfig content
notepad C:\Users\lenovo\.kube\config
\end{lstlisting}

\subsection{Step 3.9: Deploy Application to Kubernetes}

Open \textbf{PowerShell}:

\begin{lstlisting}[language=bash]
cd C:\Users\lenovo\Desktop\ass1-645

kubectl apply -f deployment.yaml

kubectl apply -f service.yaml
\end{lstlisting}

Verify deployment:

\begin{lstlisting}[language=bash]
kubectl get pods

kubectl get services
\end{lstlisting}

You should see \textbf{3 pods} running!

Copy the \textbf{EXTERNAL-IP} from the service - this is your Kubernetes URL!

\newpage

% ============================================
\section{Part 4: Set Up Jenkins}
% ============================================

\subsection{Step 4.1: Launch EC2 Instance for Jenkins}

\begin{enumerate}
    \item Go to \textbf{EC2} $\rightarrow$ \textbf{Launch Instance}
    \item Configure:
\end{enumerate}

\begin{center}
\begin{tabular}{|l|l|}
\hline
\textbf{Setting} & \textbf{Value} \\
\hline
Name & Jenkins-Server \\
AMI & Ubuntu Server 22.04 LTS \\
Instance Type & t2.medium \\
Key Pair & Create new: jenkins-key \\
\hline
\end{tabular}
\end{center}

\textbf{Security Group Rules:}

\begin{center}
\begin{tabular}{|l|l|l|}
\hline
\textbf{Type} & \textbf{Port} & \textbf{Source} \\
\hline
SSH & 22 & My IP \\
Custom TCP & 8080 & Anywhere (0.0.0.0/0) \\
\hline
\end{tabular}
\end{center}

\subsection{Step 4.2: Connect to Jenkins EC2}

\begin{lstlisting}[language=bash]
cd C:\Users\lenovo\Downloads

ssh -i "jenkins-key.pem" ubuntu@YOUR-JENKINS-PUBLIC-IP
\end{lstlisting}

\subsection{Step 4.3: Install Java}

\begin{lstlisting}[language=bash]
sudo apt update

sudo apt install -y openjdk-11-jdk
\end{lstlisting}

\subsection{Step 4.4: Install Jenkins}

\begin{lstlisting}[language=bash]
curl -fsSL https://pkg.jenkins.io/debian-stable/jenkins.io-2023.key | \
    sudo tee /usr/share/keyrings/jenkins-keyring.asc > /dev/null

echo deb [signed-by=/usr/share/keyrings/jenkins-keyring.asc] \
    https://pkg.jenkins.io/debian-stable binary/ | \
    sudo tee /etc/apt/sources.list.d/jenkins.list > /dev/null

sudo apt update

sudo apt install -y jenkins

sudo systemctl start jenkins

sudo systemctl enable jenkins
\end{lstlisting}

\subsection{Step 4.5: Install Docker on Jenkins Server}

\begin{lstlisting}[language=bash]
sudo apt install -y docker.io

sudo usermod -aG docker jenkins

sudo systemctl restart jenkins
\end{lstlisting}

\subsection{Step 4.6: Install kubectl on Jenkins Server}

\begin{lstlisting}[language=bash]
sudo snap install kubectl --classic
\end{lstlisting}

\subsection{Step 4.7: Configure kubectl for Jenkins User}

\begin{lstlisting}[language=bash]
sudo su jenkins

mkdir -p ~/.kube

nano ~/.kube/config
\end{lstlisting}

Paste the Kubeconfig content from Rancher (Step 3.8), then:
\begin{itemize}
    \item Press \texttt{Ctrl + X}
    \item Press \texttt{Y}
    \item Press \texttt{Enter}
\end{itemize}

Test the connection:

\begin{lstlisting}[language=bash]
kubectl get nodes

exit
\end{lstlisting}

\subsection{Step 4.8: Access Jenkins Web UI}

Open browser: \texttt{http://YOUR-JENKINS-PUBLIC-IP:8080}

Get the initial admin password:

\begin{lstlisting}[language=bash]
sudo cat /var/lib/jenkins/secrets/initialAdminPassword
\end{lstlisting}

Copy and paste the password in browser.

\begin{enumerate}
    \item Click \textbf{Install suggested plugins}
    \item Create admin user
    \item Save Jenkins URL
\end{enumerate}

\subsection{Step 4.9: Install Jenkins Plugins}

\begin{enumerate}
    \item Go to \textbf{Manage Jenkins} $\rightarrow$ \textbf{Manage Plugins}
    \item Click \textbf{Available} tab
    \item Search and install:
    \begin{itemize}
        \item Docker Pipeline
        \item GitHub Integration
        \item Credentials Binding
    \end{itemize}
    \item Click \textbf{Install without restart}
\end{enumerate}

\subsection{Step 4.10: Add Docker Hub Credentials}

\begin{enumerate}
    \item Go to \textbf{Manage Jenkins} $\rightarrow$ \textbf{Manage Credentials}
    \item Click \textbf{System} $\rightarrow$ \textbf{Global credentials} $\rightarrow$ \textbf{Add Credentials}
    \item Configure:
    \begin{itemize}
        \item Kind: \textbf{Username with password}
        \item Username: Your Docker Hub username
        \item Password: Your Docker Hub password
        \item ID: \texttt{dockerhub-credentials}
    \end{itemize}
    \item Click \textbf{OK}
\end{enumerate}

\subsection{Step 4.11: Create Jenkins Pipeline}

\begin{enumerate}
    \item Click \textbf{New Item}
    \item Enter name: \texttt{studentsurvey-pipeline}
    \item Select \textbf{Pipeline}
    \item Click \textbf{OK}
    \item Configure:
    \begin{itemize}
        \item \textbf{Build Triggers:} Check \textbf{Poll SCM}
        \item Schedule: \texttt{* * * * *} (every minute)
        \item \textbf{Pipeline Definition:} Pipeline script from SCM
        \item SCM: Git
        \item Repository URL: \texttt{https://github.com/YOUR\_USERNAME/studentsurvey645.git}
        \item Branch: \texttt{*/main}
        \item Script Path: \texttt{Jenkinsfile}
    \end{itemize}
    \item Click \textbf{Save}
\end{enumerate}

\subsection{Step 4.12: Run the Pipeline}

\begin{enumerate}
    \item Click \textbf{Build Now}
    \item Watch the build progress
    \item All stages should turn \textbf{green}
\end{enumerate}

\newpage

% ============================================
\section{Part 5: Testing and Verification}
% ============================================

\subsection{Step 5.1: Verify 3 Pods Running}

\begin{lstlisting}[language=bash]
kubectl get pods
\end{lstlisting}

Expected output:
\begin{lstlisting}
NAME                                        READY   STATUS    
studentsurvey-deployment-xxxxx-xxxxx        1/1     Running   
studentsurvey-deployment-xxxxx-xxxxx        1/1     Running   
studentsurvey-deployment-xxxxx-xxxxx        1/1     Running   
\end{lstlisting}

\subsection{Step 5.2: Test Pod Resiliency}

Delete one pod:

\begin{lstlisting}[language=bash]
kubectl delete pod studentsurvey-deployment-xxxxx-xxxxx
\end{lstlisting}

Check pods again - a new pod should be created automatically:

\begin{lstlisting}[language=bash]
kubectl get pods
\end{lstlisting}

\subsection{Step 5.3: Get Application URL}

\begin{lstlisting}[language=bash]
kubectl get services
\end{lstlisting}

Copy the \textbf{EXTERNAL-IP} and open in browser.

\newpage

% ============================================
\section{Part 6: Video Recording Requirements}
% ============================================

Your video must demonstrate:

\begin{enumerate}
    \item The setup process
    \item Deployment on Kubernetes cluster using kubectl
    \item The CI/CD pipeline execution in Jenkins
    \item The running application on Kubernetes
    \item The AWS URL of the deployed application
\end{enumerate}

\subsection{Recommended Recording Steps}

\begin{enumerate}
    \item Show your GitHub repository with all files
    \item Show Docker Hub with your image
    \item Show Rancher dashboard with cluster running
    \item Run \texttt{kubectl get pods} to show 3 pods
    \item Delete a pod and show it recreating
    \item Show Jenkins pipeline running
    \item Open the application URL in browser
    \item Navigate through the website
\end{enumerate}

\newpage

% ============================================
\section{Live URLs (Update After Deployment)}
% ============================================

\begin{center}
\begin{tabular}{|l|l|}
\hline
\textbf{Service} & \textbf{URL} \\
\hline
S3 Static Website & http://div-645-assignment1.s3-website.us-east-2.amazonaws.com \\
EC2 Instance (HW1) & http://3.138.154.222 \\
Kubernetes LoadBalancer & http://YOUR-K8S-LOADBALANCER-URL \\
Jenkins & http://YOUR-JENKINS-IP:8080 \\
Rancher & https://YOUR-RANCHER-IP \\
Docker Hub & https://hub.docker.com/r/YOUR\_USERNAME/studentsurvey645 \\
GitHub & https://github.com/YOUR\_USERNAME/studentsurvey645 \\
\hline
\end{tabular}
\end{center}

\newpage

% ============================================
\section{Troubleshooting}
% ============================================

\subsection{Docker Issues}

\begin{center}
\begin{tabular}{|p{5cm}|p{8cm}|}
\hline
\textbf{Problem} & \textbf{Solution} \\
\hline
Docker build fails & Check Dockerfile syntax and file paths \\
Cannot push to Docker Hub & Run \texttt{docker login} first \\
Image not found & Verify image name and tag \\
\hline
\end{tabular}
\end{center}

\subsection{Kubernetes Issues}

\begin{center}
\begin{tabular}{|p{5cm}|p{8cm}|}
\hline
\textbf{Problem} & \textbf{Solution} \\
\hline
Pods not starting & Run \texttt{kubectl describe pod <pod-name>} \\
Service not accessible & Verify security groups allow port 80 \\
kubectl not working & Check kubeconfig file path and content \\
\hline
\end{tabular}
\end{center}

\subsection{Jenkins Issues}

\begin{center}
\begin{tabular}{|p{5cm}|p{8cm}|}
\hline
\textbf{Problem} & \textbf{Solution} \\
\hline
Build fails & Check Jenkins console output for errors \\
Cannot connect to Docker & Add jenkins user to docker group \\
kubectl fails in pipeline & Verify kubeconfig in /home/jenkins/.kube/config \\
\hline
\end{tabular}
\end{center}

\subsection{Useful Commands}

\begin{lstlisting}[language=bash]
# View pod logs
kubectl logs <pod-name>

# Describe deployment
kubectl describe deployment studentsurvey-deployment

# Check service status
kubectl get svc studentsurvey-service

# Restart deployment
kubectl rollout restart deployment/studentsurvey-deployment

# Check Jenkins logs
sudo journalctl -u jenkins
\end{lstlisting}

\newpage

% ============================================
\section{Cleanup (After Grading)}
% ============================================

To avoid AWS charges after grading:

\begin{enumerate}
    \item \textbf{Rancher:} Delete cluster in Rancher UI
    \item \textbf{EC2:} Terminate all instances (Rancher-Server, Jenkins-Server, Worker nodes)
    \item \textbf{Load Balancer:} Delete in AWS EC2 $\rightarrow$ Load Balancers
    \item \textbf{Security Groups:} Delete custom security groups
    \item \textbf{Key Pairs:} Delete if no longer needed
\end{enumerate}

\newpage

% ============================================
\section{Submission Checklist}
% ============================================

\begin{itemize}
    \item[$\square$] Source code files (index.html, survey.html, error.html, campus.jpg)
    \item[$\square$] Dockerfile
    \item[$\square$] deployment.yaml
    \item[$\square$] service.yaml
    \item[$\square$] Jenkinsfile
    \item[$\square$] Documentation (PDF)
    \item[$\square$] Video demonstration with voice-over
    \item[$\square$] All URLs updated in documentation
    \item[$\square$] Name on all files
    \item[$\square$] Zipped package submitted to Canvas
\end{itemize}

\vfill

\begin{center}
\textbf{--- End of Document ---}
\end{center}

\end{document}
